\chapter{Evaluation}

The implementation was tested on all 20 \ac{BLIF} benchmark files (0--19). Table~\ref{tab:results} compares the initial cube count, the reference minimized count, and the result produced by this implementation. File~11 ($^*$) ran to completion but achieved no reduction (1027 $\to$ 1027), with \texttt{Expand} taking $\approx$45\,s and \texttt{Irredundant} $\approx$10\,min.

\begin{table}[h]
\centering
\begin{tabular}{@{}rrrr@{}}
\toprule
File & Initial & Reference & This work \\
\midrule
0  & 4   & 3   & \textbf{2}   \\
1  & 4   & 1   & 1   \\
2  & 12  & 4   & 4   \\
3  & 8   & 5   & 5   \\
4  & 14  & 7   & 7   \\
5  & 15  & 8   & 8   \\
6  & 28  & 15  & 15  \\
7  & 34  & 17  & 17  \\
8  & 69  & 29  & 29  \\
9  & 64  & 25  & 26  \\
10 & 290 & 182 & 186 \\
11 & 1027 & 771 & 1027$^*$ \\
12 & 156 & 87  & 87  \\
13 & 307 & 117 & 132 \\
14 & 240 & 10  & 10  \\
15 & 481 & 218 & 218 \\
16 & 420 & 86  & 148 \\
17 & 32  & 1   & 1   \\
18 & 46  & 46  & 46  \\
19 & 92  & 20  & 20  \\
\bottomrule
\end{tabular}
\caption{Cube counts before and after minimization.\ainote{Claude Sonnet 4.6}{Format experimental results into a LaTeX table.}}
\label{tab:results}
\end{table}

15 out of 20 tested files match or improve on the reference. File~0 produces a smaller cover than the reference (2 vs.\ 3), indicating the algorithm found a valid alternative minimization. Files~9, 10, 13, and 16 are slightly above the reference, suggesting the Reduce--Expand--Irredundant loop converged to a local minimum rather than the global optimum, which is expected behaviour for the heuristic Espresso algorithm. File~11 produced no reduction at all despite completing.

\section{Functional Equivalence}

Correctness of the minimized outputs was verified using the \texttt{cec} (combinational equivalence checking) command in the ABC logic synthesis tool. The minimized cover was checked against the original for a representative selection of files:

\begin{table}[h]
\centering
\begin{tabular}{@{}lll@{}}
\toprule
File & Result & Time \\
\midrule
1.blif  & Equivalent (structural hashing) & 0.00\,s \\
10.blif & Equivalent                      & 0.02\,s \\
13.blif & Equivalent                      & 0.03\,s \\
18.blif & Equivalent (structural hashing) & 0.00\,s \\
\bottomrule
\end{tabular}
\caption{ABC \texttt{cec} equivalence check results.\ainote{Claude Sonnet 4.6}{Format ABC CEC results into a LaTeX table.}}
\label{tab:cec}
\end{table}

All checked outputs are functionally equivalent to their inputs, confirming that the minimization preserves the Boolean function. This includes file~13, where the cube count is above the reference, verifying that the result is still correct despite being suboptimal.